\chapter{Conclusion and Future Work}\label{ch:conclusion}


This dissertation set out to show that useful near-term quantum computation is a property of a co-designed stack, not of a quantum processor alone. It also set out to make concrete progress at each layer. This chapter answers the five research questions of \cref{ch:intro}. It then draws out the design principles that recur across layers, records the remaining limitations, and lays out the research program that follows.

\section{Answers to the Research Questions}\label{sec:conclusion:answers}

\subsection{RQ1: Scale}\label{sec:conclusion:rq1}

\Cref{ch:qgear} asked how far classical HPC can push exact circuit simulation, and how to make that capability usable without rewriting programs. Technically, \qgear{} pushes exact state-vector simulation to 42 qubits on 1,024 A100 GPUs of Perlmutter with close to full GPU utilization. A laptop-scale \qiskit{} circuit reaches that scale through a constant-time gate-by-gate transformation into a \cudaq{} kernel, plus a containerized, Slurm-native deployment~\cite{guo2025qgear}.

Practically, the speed-ups that matter to a working researcher change which experiments are feasible in a day. A 34-qubit random circuit that occupies a CPU node for a day completes in about a minute on four GPUs. A 20-qubit quantum Fourier transform runs one to two orders of magnitude faster than in \qiskit{} Aer. An 18-qubit variational eigensolver loop drops from ten minutes to four seconds. \qgear{} then served as the validation and reward-evaluation engine of \cref{ch:nisq,ch:qml,ch:synthesis}, the strongest evidence that it met the usability goal.

The memory wall of exact simulation stands at roughly 42--44 qubits, even for a leadership-class allocation. Beyond that scale, simulation must become approximate (tensor networks, Clifford-plus-magic decompositions) or give way to the hardware.

\subsection{RQ2: Data}\label{sec:conclusion:rq2}

\Cref{ch:qml} asked how classical data should enter and leave a quantum processor so that learning models are shot-efficient, hardware-compatible and trainable. The four primitives developed there share one answer:
\begin{itemize}
  \item Load many values at once into an address--data register with a uniformly controlled rotation.
  \item Read many expectation values at once from the same circuit.
  \item Keep the nonlinearity on the classical side.
  \item Let a hardware-aware compiler decide how the circuit meets the device.
\end{itemize}

\qpie{} showed that parallel exchange between pretrained classical parameters and variational circuits improves accuracy over classical transfer learning on five benchmark tasks. It also speeds convergence on stochastic time series~\cite{guo2025qpie}. \vqt{} showed that a vectorized quantum dot product can estimate a transformer's attention matrix with logarithmic circuit complexity in the batch--sequence product. The estimate reached an RMSE of 0.059 on \texttt{ibm\_kingston}, and the resulting language model is competitive with prior quantum sequence models on the Brown corpus~\cite{guo2025vqt}.

\shardq{} showed that a deterministic, distance-ordered cut selection with global reconstruction lowers encoding error by about 15\% on Heron devices. It also lifts reconstruction fidelity from 0.79 to 0.90 for seven-address-qubit images~\cite{guo2025shardq}. \nnqa{} showed that a trained polynomial network compiles into exact quantum arithmetic. That arithmetic reaches above 99.5\% accuracy for polynomials up to degree 35 on both superconducting and trapped-ion hardware. Its resource footprint is an order of magnitude below variational alternatives~\cite{guo2026nnqa}.

\subsection{RQ3: Noise}\label{sec:conclusion:rq3}

\Cref{ch:nisq} asked how optimization algorithms on NISQ devices can exploit problem and hardware structure instead of black-box parameter search. \deal{} showed that a QUBO instance carries enough structure to initialize QAOA angles directly through qubit-prioritization normalization. Calibration data can drive a dynamic physical mapping, and a lightweight digital zero-noise extrapolation can be folded into the same pipeline. On IBM Heron processors, this raised the ground-state success rate of 20-qubit, 10-layer circuits by 14.81\% (\texttt{ibm\_torino}) and 7.40\% (\texttt{ibm\_marrakesh}) over vanilla QAOA. It also let seven layers do the work of nine~\cite{guo2025deal}.

\qawa{} approached the same question from the data side. Its classical-data-traceable oracle has two-qubit depth linear in the qubit count. With classical post-processing of mid-circuit measurements, it yields solutions whose quality can be verified in polynomial time. On a 156-qubit device, \qawa{} produced a portfolio 5.4\% better than the baseline and 8.1\% above the noise floor~\cite{guo2025qawa}. The general answer is to move structure out of the optimizer and into the ansatz, initialization and compilation, where it costs no shots.

\subsection{RQ4: Synthesis}\label{sec:conclusion:rq4}

\Cref{ch:synthesis} asked whether learned, HPC-scale synthesis can produce circuits that respect fault-tolerant resource budgets, and what honest cost accounting says about end-to-end quantum advantage. \rubriq{} answered the first half. It fine-tunes a code-generating language model with group relative policy optimization against a programmatic rubric, with GPU simulation inside the reward loop. The model achieves a mean T-gate compression of 3.31$\times$, against 2.05$\times$ for sparse-reward reinforcement learning. It keeps hardware-constraint violations below 1\% and converges two to three times faster. The dense, decomposable reward makes the difference~\cite{guo2026rubriq}.

The nonlinear-wave study answered the second half with a quantitative negative result. In a split-step solver that must measure the field to evaluate the nonlinearity, the per-step quantum cost (depth times shots) grows faster than the classical cost with grid size. An end-to-end advantage therefore requires a coherent nonlinear update that does not read out the field~\cite{guo2026nonlinearwaves}. Both halves rely on the same HPC substrate as \cref{ch:qgear}, and together they set the bar that fault-tolerant compilation must clear.

\subsection{RQ5: Trust}\label{sec:conclusion:rq5}

\Cref{ch:qec,ch:pqc} asked how error-corrected computation on shared cloud QPUs can be made auditable and secure, and how that connects to post-quantum cryptography standards. The audit of \cref{ch:qec} introduced a bounded cloud fault model and the notion of accepted logical disturbance, with its exact expectation under Haar-random encoders. It also introduced a polynomial-cost seeded Clifford ensemble that reseeds the encoder per run. For eight-qubit encoders, the reseeded ensemble reduces mean accepted disturbance from 0.150 to 0.020, an 86.7\% reduction attributable to syndrome-based rejection. The fixed \code{5}{1}{3} code corrects every tested weight-one fault, and gate-level Stim simulations up to seventeen qubits confirm the scaling~\cite{guo2026qecaudit}.

\Cref{ch:pqc} then showed that a post-quantum cryptographic module designed to the FIPS 140-3 requirements should supply three things. These are the seed for that encoder, the keys that protect job submission, and the signatures that authenticate results. It described the Light Rider QPC SDK as such a module, with a discovery service for approved algorithm suites~\cite{guo2026lightriderqpc}.

Trust in the quantum stack is therefore a chain. Validated randomness feeds reseeded encoders. Reseeded encoders bound what a co-tenant adversary can do. Post-quantum key establishment and signatures protect everything that moves between the classical and quantum sides.

\section{Cross-Cutting Principles}\label{sec:conclusion:principles}

Four principles recur across the chapters and are, in the candidate's view, the durable contribution of the dissertation beyond its individual results.

The first is \emph{move work off the QPU and out of the shot budget}:
\begin{itemize}
  \item \deal{} moves parameter search into initialization and compilation.
  \item \qawa{} and \vqt{} move nonlinear post-processing to the classical side.
  \item \nnqa{} moves training entirely off the device and uses it only for exact arithmetic.
  \item \rubriq{} moves the search for good circuits into a classical policy trained against a simulator.
\end{itemize}
In every case, the quantum processor does only what the classical side cannot.

The second is \emph{vectorize the interface}. Address--data encodings and parallel readout appear in four chapters:
\begin{itemize}
  \item in \cref{ch:qgear} as the \qcrank{} benchmark;
  \item in \cref{ch:qml} as the foundation of all four primitives;
  \item in \cref{ch:nisq} as the data-traceable oracle; and
  \item in \cref{ch:synthesis} as the reload step of the split-step solver.
\end{itemize}
Whatever the application, the number of values that cross the classical--quantum interface per circuit execution sets its cost, and vectorization is the lever.

The third is \emph{let calibration data drive compilation}. \deal's dynamic mapping and \rubriq's hardware-compliance term take device calibration and topology as compilation inputs, not as an afterthought. \shardq{} transpiles its sub-circuits against the target's calibrated coupling map. The devices used here changed their error maps daily, and methods that read the map outperformed those that did not.

The fourth is \emph{account honestly and audit adversarially}. The nonlinear-wave study reports where the quantum cost overtakes the classical one. The QEC audit assumes an adversary rather than benign noise, and \cref{ch:pqc} distinguishes what is validated from what is designed. A stack trusted with computations of value must be assessed against the costs it will actually incur and the adversaries it will actually face.

\section{Limitations}\label{sec:conclusion:limitations}

The limitations recorded in each chapter fall into three groups. The first is scale on the quantum side. Hardware experiments used at most 36 qubits and a few hundred two-qubit gates. Those experiments showed that the algorithms of \cref{ch:nisq,ch:qml} help at that scale, and simulation extends the evidence to 42 qubits. Whether the observed gains persist at hundreds of logical qubits is an open empirical question.

The second is model scope. The QEC audit uses weight-one Pauli faults and small codes. The cost model of the nonlinear-wave study assumes a particular loader and polynomial nonlinearity. \rubriq's rubric encodes one view of fault-tolerant cost centred on T-count. Each model is tractable and defensible by design, and each could be broadened.

The third is validation status. Several results are preprints under review at the time of writing. The cryptographic module of \cref{ch:pqc} is designed to a standard rather than certified against it. Each chapter states these facts where they arise, and the artifact appendix (\cref{app:repro}) records the status of each work.

\section{Future Work}\label{sec:conclusion:future}

The research program that follows has one item per layer and one that spans them.

At the \emph{scale} layer, the natural extension of \qgear{} is approximate simulation beyond the memory wall. This means matrix-product-state and tree-tensor-network backends with bond dimension chosen automatically from the circuit's entanglement structure. It also means Clifford-plus-magic decompositions that exploit the low T-count of the circuits \rubriq{} produces. A second extension is hardware-in-the-loop simulation, which refreshes the noise model from the device's live calibration. The simulator then predicts hardware results rather than ideal ones.

At the \emph{data} layer, the four primitives of \cref{ch:qml} should be unified into one vectorized-encoding library with a common compiler pass built on \shardq{}. \vqt{} should be scaled to sequence lengths and vocabularies where the vectorized dot product's logarithmic circuit complexity becomes the dominant advantage. \nnqa's exact arithmetic suggests a broader program: compile classically trained models into quantum arithmetic blocks. These include the coherent nonlinear update that \cref{ch:synthesis} identifies as the missing piece for quantum simulation of nonlinear dynamics.

At the \emph{algorithms} layer, \deal{} and \qawa{} should be combined. The direct initialization and dynamic mapping of \deal{} would join the data-traceable oracle and verifiable post-processing of \qawa{}, applied to constrained problems with penalty structure. The longer-term goal is their integration with error-corrected execution, in which the logical-qubit error map replaces the physical one.

At the \emph{synthesis} layer, \rubriq's rubric should be extended from T-count to full fault-tolerant resource estimates (logical qubits, code distance, magic-state factories). The policy then optimizes what a fault-tolerant machine will actually charge for. For larger circuits, the approximate backends proposed above should replace the simulator in the reward loop.

At the \emph{trust} layer, the QEC audit should be extended to correlated and higher-weight faults, to surface and bivariate-bicycle codes, and to an adversary who observes syndromes. The QPC SDK should proceed through formal FIPS 140-3 validation, with the hybrid key-establishment and signed-result protocols of \cref{ch:pqc} exercised against live cloud QPU services.

Spanning all layers, the most important next step is an end-to-end demonstration of one workload:
\begin{itemize}
  \item developed in simulation with \qgear;
  \item compiled with \shardq{} and \rubriq;
  \item executed with \deal-style initialization on a cloud QPU under a reseeded error-correcting encoder whose seeds and job channels are protected by the QPC SDK; and
  \item with the full cost accounted for.
\end{itemize}
That demonstration would test the thesis of this dissertation directly, and the components to build it now exist.

\section{Concluding Remarks}\label{sec:conclusion:remarks}

The quantum processor is the most photographed part of a quantum computer, but it is not where most of the computation happens. It will not be for some time. This dissertation treats the classical HPC around the QPU as a first-class object of research. It has argued, and shown across simulation, algorithms, data encoding, synthesis and security, that this treatment produces gains at every layer. Those gains compound. The stack is the machine.
